\begin{frame}[fragile]{The freeze instruction}
The \texttt{freeze} instruction is a \textbf{no-op}, unless its input is poison, in which case it non-deterministically chooses an \textbf{arbitrary} legal value.

\begin{lstlisting}[style=llvmmono]
    %x = freeze(poison) ; (*@x\leftarrow 42@*)
    %y = freeze(poison) ; (*@y\leftarrow 23@*)
\end{lstlisting}
\vspace{1.5em}
The semantics of the \texttt{r = freeze i<sz> op}
\begin{columns}[T]
    \begin{column}{0.44\textwidth}
         \begin{prooftree}
           \AxiomC{$\mathscr{E}\llbracket op \rrbracket_R = v \neq \textbf{poison}$}
            \UnaryInfC{$\text{R,M} \hookrightarrow \text{R}[r \mapsto v]\text{,M}$}
        \end{prooftree}       \begin{center}
        \end{center}
    \end{column}
    \begin{column}{0.44\textwidth}
        \begin{center}
        \begin{prooftree}
          \AxiomC{$\mathscr{E}\llbracket op \rrbracket_R = \textbf{poison} \quad v \in Num(sz)$}
            \UnaryInfC{$\text{R,M} \hookrightarrow \text{R}[r \mapsto v]\text{,M}$}
        \end{prooftree}
        \end{center}
    \end{column}
\end{columns}

\pdfpcnote{
Freeze is a no-op on normal values and a non-deterministic concretization on
poison. Crucially, each invocation is independent — the example shows two freezes
of the same poison yielding different values (42 and 49). The formal rules on the
right show the two cases: pass-through for normal values, arbitrary choice from
Num(sz) for poison. **Note:** that freeze is deterministic per execution but
non-deterministic across executions.
}

\end{frame}
